\section{Discussion and Open Problems}\label{sec:discussion}

The counterexample is governed by a simple local event that becomes visible
only after squaring the operator. A negative quadrilateral flux removes its
odd-displacement couplings from \(A^{2}\), whereas a positive flux activates them
with amplitude two. The all-negative phase is therefore a natural
cancellation baseline at squared radius eight. A pair of positive defects can
remain below that barrier only when it is antipodal in an eight-site cell; all
other period-eight geometries create enough signed closed-walk mass to cross
the barrier.

This mechanism explains why the conjecture survives many small finite orders.
The target phase has an infinite-volume edge slightly below eight, while the
twisted conjectured threshold increases toward eight. Their crossing is first
realized at order 32. The same period-eight word then works at every larger
multiple of eight because its Floquet spectrum is controlled uniformly over
the full unit circle.

The arbitrary-period moment identities show that the defect picture is not an
artifact of period eight. Any phase at or below the eight-barrier must have
defect density at most \(\frac{3}{4}\), and its adjacent and distance-two defect clusters
must satisfy a stronger weighted inequality. These constraints are only
necessary. The low-period frontier illustrates why: sparse near-barrier
geometries can require closed walks of length up to 130 before a positive
excess appears, and five bounded-period competitors are more efficiently
excluded by endpoint Rayleigh vectors.

The computer-assisted results have two different functions. The finite search
through order 30 establishes minimality of the counterexample. The bounded
periodic classification establishes a structural frontier around the target.
Neither computation proves an unbounded optimization theorem.

\subsection{Controlled evidence beyond period sixteen}

To probe the boundary of the periodic classification, we performed a separate
exploratory calculation for displayed periods \(17\leq p\leq24\). This
calculation is not part of any theorem. It enumerates all legal \(Q\)-orbits
exactly and applies the first two moment excesses to every orbit. An independent
Burnside count and a second record-level orbit partition check completeness.
The exact orbit counts increase from 2,056 at period 17 to 176,906 at period 24. At
period 24, for example, \(F_1>0\) excludes 269 orbits and \(F_2>0\), after the
first filter, excludes another 131,290; 45,347 orbits remain unclassified by
these two implications.

We then evaluated a declared deterministic selection of 256 residual candidates
per period on finite Bloch grids and refined the best sixteen. These values are
floating exploratory data, not certified continuous-fiber bounds. No selected
candidate below \(\eta\) appeared. At period 24 the best row is exactly the
threefold displayed-cell repetition of the target, recognized as primitive
period eight; its value \(R=\eta\) follows from zone folding, not from the grid.
The next sampled period-24 candidate has primitive \(\tau\)-period 24 and grid
value approximately \(7.974985\). The complete low-moment residual sets were
not spectrally excluded, so this evidence neither extends
Theorem~\ref{thm:low-period-frontier} nor establishes an all-period conjecture.

The current public-status check found no direct public resolution in the
source paper, the author's listed follow-up materials, or the narrow direct
searches performed on 21 August 2026. That assessment should be refreshed
before submission and should not be read as an absolute priority statement.

The work leaves several natural questions.

\textbf{Problem 1.} Is the target phase globally optimal among periodic phases of
arbitrary primitive period?

\textbf{Problem 2.} Is there a meaningful infinite-volume class containing
nonperiodic signings in which the target remains optimal?

\textbf{Problem 3.} Which even orders not divisible by eight admit counterexamples,
and what is the exact finite minimum at those orders?

\textbf{Problem 4.} Can the hierarchy \(F_k=M_{k+1}-8M_k\) be converted into a finite
structural classification at arbitrary period, or do genuinely new local
motifs appear at every scale?

\textbf{Problem 5.} Is there a flux-phase variational principle that explains the
antipodal defect geometry without finite orbit classification?

Answers to these questions would move the present bounded and necessary
results toward an unbounded flux-phase theorem. They are not assumed anywhere
in this paper.
